\documentclass[10pt,a4paper]{amsart}
\usepackage[margin=19mm]{geometry}
\usepackage{amsmath,amssymb,mathtools}
\usepackage[scr=rsfs,cal=euler]{mathalfa}
\usepackage{xfrac}
\usepackage[unicode,hidelinks,bookmarks=false]{hyperref}
\newcommand{\ie}{i.e.\ }
\newcommand{\cf}{cf.\ }
\newcommand{\resp}{respectively\ }
\newcommand{\Subsection}{\subsection}
\providecommand{\nobreakdash}{\nobreak\hbox{-}}
\setlength{\emergencystretch}{2em}
\multlinegap0pt
\usepackage{amssymb}
\usepackage{booktabs}
\usepackage{mathtools}
\usepackage{bm}
\usepackage{caption}
\usepackage{xcolor}
\usepackage{tikz}
\newtheorem{proposition}{Proposition}
\title{Equivariant Cohomology, BRST Quantization, and Analytic Localization: A Unified Framework}
\author{Lixin Xu}
\date{Source: arXiv:2601.00256v1 (CC BY 4.0)}
\begin{document}
\maketitle

\noindent\emph{Source and licence.} Equivariant Cohomology, BRST Quantization, and Analytic Localization: A Unified Framework. Version arXiv:2601.00256v1 (CC BY 4.0), distributed by arXiv under the Creative Commons Attribution 4.0 International licence, http://creativecommons.org/licenses/by/4.0/, as recorded in the arXiv licence metadata for this exact version and shown on the abstract page https://arxiv.org/abs/2601.00256v1. Full licence notice: \texttt{licenses/arxiv-2601.00256-CC-BY-4.0.txt}.\par\smallskip

\noindent\textit{Editorial scope.} Counted unit: the article's proposition proposition (source0.tex lines 340--342). The article's complete original proof of it is delivered with it (source0.tex lines 343--355, one \texttt{proof} environment of 13 lines). No mathematics is written, repaired or re-derived: the statement and the proof are the article's own lines, in the article's own notation, and no retained line is substituted or re-flowed.  No further source-local context is needed: every cross-reference of every retained line resolves inside the two delivered spans.  Nothing was omitted from any retained span, so the package carries no omission record.  No retained line contains a citation, so the wrapper carries no bibliography and no bibliography entry.  No retained line contains a figure environment, an image-inclusion command or drawing code, so no image is delivered and the package declares no asset dependency.  Editorial formatting only, in addition: an AMS \texttt{amsart} preamble that restates the article's own declarations and macros used by the retained lines, namely the environments \texttt{proposition} on separate counters, together with the standard packages those lines need.  The retained environments print in this document as Proposition 1 in order of appearance, whereas the article prints the same items with the numbering of its own full document.  The complete native source of the article as distributed in the arXiv e-print for this exact version is bundled unchanged as \texttt{sources/191984/source0.tex}.
\begin{proposition}[$t$-independence]
$I(t)$ is independent of $t$: $I(t) = I(0)$ for all $t \ge 0$.
\end{proposition}

\begin{proof}
The map
\begin{equation}
\Phi_t: \alpha \mapsto e^{-tf} \alpha,
\end{equation}
is an isomorphism of complexes from $(\Omega_G^\bullet(M), d_G)$ to $(\Omega_G^\bullet(M), d_{G,t})$, because $d_{G,t} = e^{-tf} d_G e^{tf}$. Specifically, $\Phi_t$ is a chain map because
\begin{equation}
d_{G,t} (\Phi_t \alpha) = (d_G + t \, df \wedge) (e^{-tf} \alpha) = e^{-tf} d_G \alpha - t (df \wedge e^{-tf} \alpha) + t (df \wedge e^{-tf} \alpha) = e^{-tf} d_G \alpha = \Phi_t (d_G \alpha),
\end{equation}
where we used that $d_G (e^{-tf} \alpha) = e^{-tf} d_G \alpha + (d_G e^{-tf}) \alpha and d_G e^{-tf} = -t df e^{-tf}$. The map is invertible, with inverse $\Phi_t^{-1} \alpha = e^{tf} \alpha$.

Therefore, it induces an isomorphism on equivariant cohomology. In particular, the cohomology class of $e^{-tf} \omega$ in the $d_{G,t}$-cohomology is the image of the fixed class $[\omega]_{d_G}$. The integral $I(t) = \int_M [e^{-tf} \omega]_n$ depends only on this cohomology class, as the integral of the top-degree component of an equivariantly closed form is a well-defined functional on equivariant cohomology (since equivariantly exact forms integrate to zero over compact $M$ without boundary). Hence, $I(t)$ is constant in $t$, $I(t) \equiv I(0) = \int_M [\omega]_n$.
\end{proof}

\end{document}
